\documentclass[10pt]{article}
\usepackage[margin=20mm]{geometry}
\usepackage{amsmath,amssymb,mathtools,fontspec,unicode-math,xurl,hyperref}
\IfFontExistsTF{DejaVu Serif}{\setmainfont{DejaVu Serif}}{\setmainfont{Latin Modern Roman}}
\setmathfont{latinmodern-math.otf}
\setlength{\parindent}{0pt}\setlength{\parskip}{6pt}\sloppy
\title{Scientific review of
Restarts Subject to Approximate Sharpness:
A Parameter-Free and Optimal Scheme
For First-Order Methods}
\author{Math2025 scientific review — M25-0439}
\date{29 October 2026}
\begin{document}
\maketitle
Abstract: This initial scientific continuation examines the acquired Springer version of record of the paper by Adcock, Colbrook and Neyra-Nesterenko. Four substantive failures of stated guarantees are established by explicit admissible examples: the arbitrary scheduling set in Theorem 3.3 need not support its counting argument; the conditional inner-work contract does not control incorrectly guessed branches; Algorithm 6 returns a query iterate for which its claimed accuracy bound fails; and Proposition 5.2 fixes a feasibility penalty that is insufficient under its general robust null-space premise. A further local error concerns the inference of individual feasibility from a combined objective–residual bound. The report gives sufficient corrections where proved, identifies the surviving complete-grid arguments and positive checks, and distinguishes the literal claims from the reported numerical experiments. It does not claim that every method, experiment, or intended corrected theorem fails. Scientific adjudication remains pending; no independent acceptance is asserted.\par
\section*{Source identity and scope}
The reviewed work is B. Adcock, M. J. Colbrook and M. Neyra-Nesterenko, Restarts Subject to Approximate Sharpness: A Parameter-Free and Optimal Scheme For First-Order Methods, DOI 10.1007/s10208-024-09673-8 (https://doi.org/10.1007/s10208-024-09673-8) [paper]. The acquired PDF is the Springer version of record, published online on 18 February 2025 and assigned to Foundations of Computational Mathematics 26 (2026), 1137–1192. The 2025 task membership therefore refers to its online publication year. Its 56 physical PDF pages are numbered here as \(p=1,\ldots,56\); the corresponding printed page is \(1136+p\).\par
The original PDF has 3,264,645 bytes and SHA256\par
`fd6a4ef0307437644d5fc55f3cccd5ae54c97913f594df302d8bba34ff03c98c`.\par
The source manifest has 85,588 bytes and SHA256\par
`ccc3a2aafb27a8e804d31e01c3df83a071f0d013aff2b1cb14f4caca00f80d80`.\par
The related arXiv manuscript is 2301.02268v2, dated 22 July 2024; no equivalence with that earlier manuscript is needed or inferred. The retained publisher/arXiv/Crossref metadata did not expose a correction candidate in the checked records. That bounded search does not prove that no correction exists elsewhere.\par
This report continues the missing initial work for M25-0439. The recoverable material comprised source bytes, metadata, a released-code snapshot, extraction/indexing, render payloads and prepared dependencies. No usable substantive predecessor report or proof was recovered from the bounded materialized search. Those mechanical outputs were retained and reused, with new actual scientific reading and image observations distinguished from the earlier executions. This is not a second independent acceptance review.\par
The analysis covers the combined-error framework, known-parameter restart proof, the general scheduling theorem and its concrete corollaries, concrete first-order method specializations, and the applications and evidence described below. No author code or published numerical experiment was run.\par
\section*{Mathematical setup and dependency map}
Equation reference: eq:sharp, eq:inner. Write \[\mathcal E(x)=f(x)-\widehat f+g_Q(x),\qquad
 \widehat f=\min_{x\in Q}f(x),\qquad \widehat X=\operatorname*{argmin}_{x\in Q}f(x).\] The setup on pp. 2–3 assumes a proper closed convex objective with effective domain \(D\), a closed convex set \(Q\subseteq D\), a nonempty minimizer set, and a nonnegative feasibility residual \(g_Q\) that vanishes exactly on \(Q\). It also assumes that \(g_Q(x_n)\to0\) when the distance from \(x_n\) to \(Q\) tends to zero. Approximate sharpness is \[
 d(x,\widehat X)\le
 \left(\frac{\mathcal E(x)+\eta}{\alpha}\right)^{1/\beta},
 \qquad \alpha>0,\quad\beta\ge1,\quad\eta\ge0.\] The inner solver satisfies \[
 d(x_0,\widehat X)\le\delta
 \ \Longrightarrow\
 \mathcal E\bigl(\Gamma(\delta,\varepsilon,x_0)\bigr)\le\varepsilon .\] The printed iteration bound \(C_\Gamma(\delta,\varepsilon)\) is likewise specified for starts within distance \(\delta\), a distinction relevant to Finding F2.\par
Algorithm 1 uses known \(\alpha,\beta\) and unknown \(\eta\). Algorithm 2 tries a grid \(\alpha_i=a^i\alpha_0\), \(\beta_j=b^j\beta_0\), interleaving the choices through a schedule \(\phi\). Each pair has a restart count \(U_{i,j}\) and an accumulated declared-work counter \(V_{i,j}\). At a visit \((i,j,k)\), it runs at most one new restart if \(V_{i,j}+C_\Gamma\le k\), then retains the better of the old and new points by \(f+g_Q\).\par
Theorem 3.3 supplies the general work bound. Corollaries 3.4–3.7 specialize its scheduling set; Theorem 1.1 presents the resulting rates. Section 4 verifies proposed inner solvers, including the Hölder-smooth specialization in Proposition 4.10. Proposition 5.2 asserts approximate sharpness for quadratically constrained basis pursuit (QCBP), supplying a theoretical basis for one of the applications. These dependencies determine the severity assessments; no target fraction of Major findings is used.\par
\section*{Established findings}
\subsection*{F1: arbitrary scheduling sets do not justify the cardinality bound}
Location and status. Theorem 3.3, pp. 18–20, especially (3.9)–(3.10); proof p. 21. Major: its general arbitrary-set guarantee is false.\par
Equation reference: eq:sublevel. The theorem permits any infinite \(S\subseteq\mathbb Z\times\mathbb N_0\times\mathbb N\), subject to an \(h\)-assignment and membership of the target triple \((I,J,K(\varepsilon))\). It bounds the inner work by \[
 \#\{(i,j,k)\in S:
 h(|i|,j,k)\le h(|I|,J,K(\varepsilon))\}.\] Its proof assumes that each fixed pair visits \(k=1,2,\ldots\) successively and that every smaller counter value is present. Neither assertion follows from an arbitrary infinite \(S\).\par
Here is a counterexample independent of ties, incorrect runtime estimates, or positive noise floors. Take \[D=Q=\mathbb R,\quad f(x)=|x|,\quad g_Q=0,\quad
 \widehat X=\{0\},\quad \widehat f=0,\quad
 \alpha=\beta=1,\quad\eta=0.\] If necessary, regard \(\mathbb R\) as a closed real subspace of the complex ambient space and extend the objective by \(+\infty\) off it. All domain and convexity conditions hold, and [eq:sharp] (\#eq:sharp) is an equality. Define \(\Gamma(\delta,\varepsilon,x)=\varepsilon\), with one explicit update. Its error is exactly \(\varepsilon\), for every start, and its actual cost is uniformly \(C_\Gamma=1\). This cost is monotone as required and satisfies the optional polynomial-plus-one bound.\par
Set \(x^{(0)}=\varepsilon_0=1\), \(r=1/2\), \(a=b=2\), \(\alpha_0=\beta_0=1\), and target \(\varepsilon=1/4\). Then \(I=J=0\), two reductions are required, and \(K(\varepsilon)=2\). Choose \[S=\{(0,0,2n):n\in\mathbb N\},\qquad h(x_1,x_2,k)=k,\qquad
 \phi(n)=(0,0,2n).\] This \(h\) is nondecreasing in its first two arguments and strictly increasing in its third. The assignment is bijective and satisfies even the printed non-strict ordering biconditional, without ties. The required target \((0,0,2)\) lies in \(S\), but its sublevel has cardinality one.\par
The first visit has \(k=2\). The gate permits one call at radius \(2\) and accuracy \(1/2\), returning \(1/2\). The target \(1/4\) is not reached. The next visit, \(k=4\), permits a second call at radius \(1\) and accuracy \(1/4\), which reaches the target. Thus the necessary work is two, contradicting the stated upper bound one. The exact fractional trace is retained in the schedule auxiliary’s calculation record.\par
Equation reference: eq:down. A sufficient structural correction is \[
 (i,j,k)\in S\Longrightarrow(i,j,\ell)\in S
 \quad(1\le\ell\le k).\] The concrete grids in Corollaries 3.4–3.7 have this property. Consequently this counterexample does not refute their counting estimates or the corresponding concrete-grid rates after the other premises are repaired. Section 4.1 (\#sec:countrepair) proves the corrected charging argument.\par
\subsection*{F2: valid-start cost estimates do not cap incorrect branches}
Location and status. The conditional definition of \(C_\Gamma\), p. 3; Algorithm 2, p. 16, lines 11–14; explanation p. 17 and work accounting p. 21. Major: the general actual-work conclusion needs an additional premise.\par
The printed pre-call gate checks a declared cost. It does not truncate a call whose guessed radius is too small. A total solver may obey every stated valid-start guarantee and nevertheless take arbitrarily many iterations outside that radius.\par
Use \[D=\mathbb R,\quad Q=[0,\infty),\quad f(x)=|x|/8,\quad
 g_Q(x)=\max\{-x,0\},\quad d(x,\widehat X)=|x|.\] Then \(\widehat X=\{0\}\), \(\widehat f=0\), and all objective, domain and residual hypotheses hold. The combined error is \(|x|/8\) for \(x\ge0\) and \(9|x|/8\) for \(x<0\), so [eq:sharp] (\#eq:sharp) holds globally with \(\alpha=1/8,\beta=1,\eta=0\).\par
Define a terminating iterative solver on every input. If \(|x_0|\le\delta\), it takes one update and returns \(\min\{|x_0|,8\varepsilon\}\). Otherwise it takes seven unchanged iterates and returns \(x_0\). The unchanged iterates can each evaluate the objective under an oracle-work interpretation. On every valid start the output has error at most \(\varepsilon\) and the actual cost is one. Thus \(C_\Gamma=1\) meets the literal conditional contract, all monotonicities, and the polynomial-plus-one bound. No nontermination is used.\par
Set \[x^{(0)}=1,\quad \varepsilon_0=1/8,\quad
 r=1/2,\quad a=b=2,\quad \alpha_0=\beta_0=1.\] For target \(\varepsilon=1/16\), the correct pair is \(I=-3,J=0\), with \(K=1\). Let \[S=\{0,-3\}\times\{0\}\times\mathbb N,\qquad
 h(x_1,x_2,k)=2^{x_1/2}k.\] Both vertical fibers are complete, so F1 is absent. Their priorities are \(k\) and \(2\sqrt2\,k\), with no ties and finite sublevels. Sorting these values gives a bijection satisfying the printed definition. The target sublevel consists of exactly three visits: \[(0,0,1),\quad (0,0,2),\quad (-3,0,1).\]\par
Visit | Radius | Requested error | Actual work | Output\par
\((0,0,1)\) | \(1/4\) | \(1/16\) | \(7\) | \(1\)\par
\((0,0,2)\) | \(1/8\) | \(1/32\) | \(7\) | \(1\)\par
\((-3,0,1)\) | \(2\) | \(1/16\) | \(1\) | \(1/2\)\par
Every gate is satisfied with declared cost one. The first two inputs violate their radius premise; the third is valid and reaches combined error \(1/16\). The actual work is \(15>3\). Replacing seven by arbitrary positive \(N\) gives \(2N+1\) against the same claimed bound three. The saved independent calculation traces these exact radii, gates, outputs and charges; it is an arithmetic simulation of the defined algorithm, not a timing benchmark.\par
A sufficient repair is an all-start work cap: actual work must be at most \(C_\Gamma(\delta,\varepsilon)\) for every \(x_0\in D\), while accuracy remains conditional on the radius. Alternatively, for an interruptible iterative solver, explicitly stop a call after its integer budget and retain the old point if the call has not completed. Every valid call still succeeds by the original contract, and incorrect branches cannot overspend. The capped trace of this example has work three.\par
This identifies a gap in the general solver contract, not a demonstrated runtime failure of every fixed-iteration method in Section 4. Such implementations may already have the stronger all-start cap. The correct-pair error induction and the declared-budget arithmetic survive.\par
\subsection*{F3: Algorithm 6 returns the wrong sequence for its stated guarantee}
Location and status. Algorithm 6, p. 31; Lemma 4.9 and Proposition 4.10, p. 32; Hölder row of Table 1, p. 7. Major: the displayed inner-accuracy promise is false for the returned iterate.\par
In Algorithm 6, \(x_{k+1}\) is a gradient-query point, whereas \(y_{k+1}\) incorporates the proximal update. The output is explicitly \(x_N\), and Lemma 4.9 states its bound for \(x_k\).\par
Take the finite-valued convex functions \[Q=\mathbb R,\qquad q(x)=x^2/2,\qquad g(x)=10|x|,
 \qquad f=q+g.\] The minimizer is zero. The smoothness parameters are \(\nu=M_\nu=1\); the proximal map of \(g\) is elementary. Choose \[x_0=y_0=1,\quad \delta=1,\quad L_0=1,\quad
 \varepsilon=8,\quad A_0=0.\] Proposition 4.10 prescribes \[N=\left\lceil
  \frac{2^{(2+4\nu)/(1+3\nu)}
         M_\nu^{2/(1+3\nu)}
         \delta^{(2+2\nu)/(1+3\nu)}}
       {\varepsilon^{2/(1+3\nu)}}\right\rceil=1.\] At \(k=0\), the first line-search trial has \(a_1=1\) and \(\tau_1=1\). The query point is \(x_1=v_0=x_0=1\). Its proximal successor is \[y_1=\operatorname*{argmin}_{z\in\mathbb R}
 \left\{z+10|z|+\tfrac12(z-1)^2\right\}=0 .\] The backtracking acceptance inequality reads \(0\le4\), so this trial is accepted. Yet the stated output is \(x_N=x_1=1\), and \[f(x_N)-\widehat f=\tfrac{21}{2}>8=\varepsilon.\] Lemma 4.9 at \(k=1\) has right-hand side \(\sqrt{64}/2+8/2=8\) and fails for the same query point. The displayed loop includes \(k=N\), but computing later stored iterates does not change \(x_1\); that indexing ambiguity cannot rescue this example.\par
The exact first-step calculation confirms the accepted test, the two sequences and the failed bound. It executes independently written arithmetic, not an author implementation. The example uses neither an extended-valued proximal pathology nor a failed line search.\par
Returning the proximal sequence \(y_N\), with corresponding changes to the lemma and iteration indexing, is a necessary correction and fixes this example. A complete repaired universal theorem still requires verification of the cited method’s initialization and line-search hypotheses, including the treatment of arbitrary \(L_0>0\). That complete restoration is not certified here. The printed Proposition 4.10 therefore does not currently justify the Hölder row through the displayed output. The inspected experimental snapshot uses other methods; no failure of its figures follows from this example.\par
\subsection*{F4: the fixed QCBP penalty does not follow from general rNSP}
Location and status. Definition 5.1 and Proposition 5.2, pp. 38–39; introductory QCBP claim (1.6), p. 5; application parameter choice, p. 39. Major: the asserted application-level sharpness implication is false.\par
For QCBP the paper takes \[f(z)=\left\lVert z\right\rVert_1,\qquad
 Q=\{z:\left\lVert Az-y\right\rVert_2\le\zeta\},\qquad
 g_Q(z)=\sqrt{s}\max\{\left\lVert Az-y\right\rVert_2-\zeta,0\}.\] Proposition 5.2 claims [eq:sharp] (\#eq:sharp) with \[\alpha=c_1\sqrt{s},\quad \beta=1,\quad
 \eta=c_2\sigma_s(x)_1+c_3\zeta\sqrt{s},\] where positive \(c_1,c_2,c_3\) depend only on the robust null-space constants \(0<\rho<1\), \(\gamma>0\). The proposition requires \(\zeta>0\).\par
Fix \[n=m=s=1,\quad A=[1/2],\quad \rho=1/2,\quad\gamma=2,
 \quad x_{\rm true}=1,\quad y=1/2,\quad e=0.\] The rNSP holds for every real or complex \(v\): on the full support, \(\gamma|Av|=|v|\), and the empty support is trivial. The ground truth is exactly sparse, so \(\sigma_s(x_{\rm true})_1=0\).\par
Equation reference: eq:negative. For any proposed constant \(c_3>0\), set \[\zeta=\frac1{4(c_3+1)}>0.\] The feasible set is the disk with center one and radius \(2\zeta<1/2\); its unique minimum-modulus point is \(\widehat x=1-2\zeta\), with \(\widehat f=1-2\zeta\). At \(z=0\), \[f(0)-\widehat f=-1+2\zeta,\qquad
 g_Q(0)=\tfrac12-\zeta,\qquad
 \eta=c_3\zeta.\] Consequently \[
 \mathcal E(0)+\eta=-\tfrac12+(c_3+1)\zeta=-\tfrac14.\] The left side of [eq:sharp] (\#eq:sharp) is \(1-2\zeta>0\), whereas its right side for \(\beta=1\) is negative for every positive \(\alpha\). No choice of \(c_1\) or \(c_2\) repairs this instance. The constants cannot depend on the noise radius, since the stated dependence is only on the fixed \(\rho,\gamma\). Thus this family refutes the quantified proposition with its required positive radius; a zero-noise limiting argument alone would not have been sufficient.\par
The issue is the relative coefficient of the residual, which changes when \(A\) is rescaled. Section 4.2 (\#sec:qrepair) proves a conservative rNSP-dependent replacement. That replacement changes an input to the inner method and depends on \(\gamma\); it does not by itself restore a wholly parameter-free implementation. In particular, the experiment’s choice \(\kappa=\sqrt m\), justified there by \(s\le m\), is not automatically sufficient under the proposition’s general premise.\par
The specific experiments use normalized random Gaussian matrices. The scalar instance above does not demonstrate that those random draws or their displayed reconstructions failed. It invalidates the general claimed implication, not the reported data.\par
\subsection*{F5: a small combined error need not mean epsilon feasibility}
Location and status. Interpretation of (1.3), p. 3; individual-rate assertion preceding Theorem 2.1, p. 12. Minor: a local quantitative interpretation fails, while the formal combined-error conclusions survive.\par
Let \[D=[-1,1],\quad Q=[0,1],\quad f(x)=x\ \ (x\in D),\] with \(f=+\infty\) outside \(D\), and set \[g_Q(x)=
 \begin{cases}x^2-x,&x<0,\\0,&x\ge0.\end{cases}\] This residual is nonnegative and continuous, vanishes exactly on \(Q\), and satisfies the required distance-to-\(Q\) condition. We have \(\widehat X=\{0\}\), \(\widehat f=0\), and \(\mathcal E(x)=x^2\) for \(x<0\), \(\mathcal E(x)=x\) for \(x\ge0\). Thus \(|x|^2\le\mathcal E(x)\) on \(D\), proving approximate sharpness with \(\alpha=1,\beta=2,\eta=0\).\par
The uniformly one-step solver \(\Gamma(\delta,\varepsilon,x)=-\min\{\sqrt\varepsilon,1\}\) satisfies [eq:inner] (\#eq:inner) for all starts. For \(0<\varepsilon\le1\), however, \[\mathcal E(\Gamma)=\varepsilon,\qquad
 f(\Gamma)-\widehat f=-\sqrt\varepsilon,\qquad
 g_Q(\Gamma)=\sqrt\varepsilon+\varepsilon>\varepsilon .\] Algorithm 1 with \(x_0=\varepsilon_0=1\), \(r=1/4\), follows these outputs. Hence \(g_Q/\varepsilon=1+\varepsilon^{-1/2}\) is unbounded, not merely a missing constant. At \(\varepsilon=1/16\), the feasibility residual is \(5/16\). Eight exact fractional instances support the general analytic formula.\par
Always \(f-\widehat f\le\mathcal E\), because \(g_Q\ge0\). The additional conclusion \(g_Q\le\mathcal E\) requires \(f\ge\widehat f\) at the returned point. More generally, \(g_Q\ge\kappa(\widehat f-f)_+\), \(\kappa>1\), yields \[g_Q\le\frac{\kappa}{\kappa-1}\mathcal E.\] Neither extra premise is part of the general setup. When \(\eta=0\), \(\mathcal E\to0\) still implies distance to \(\widehat X\) and then to \(Q\) tends to zero, so the stipulated continuity gives qualitative feasibility. The example does not refute that qualitative conclusion or the combined-error theorems.\par
\section*{Supported corrections and surviving arguments}
\subsection*{A corrected scheduling theorem}
Impose downward completeness [eq:down] (\#eq:down), positive integer charged costs, and an all-start actual-work cap as in F2. Order the schedule by nondecreasing \(h\), with arbitrary ordering of equal priorities.\par
For a fixed pair, let \(c_q\) be the declared cost of its \(q\)-th restart and \(T_q=\sum_{p=1}^q c_p\). These costs depend on the fixed pair and its restart index, not on the interleaving. Because \(h\) is strictly increasing in \(k\), downward completeness makes the visited counter values exactly \(1,2,\ldots\). The first restart occurs at \(k=T_1\). Inductively, after restart \(q-1\), the counter is \(V=T_{q-1}\); the next gate first opens at \(k=T_q\). Since \(c_q\ge1\), one call per visit suffices. The correct pair therefore completes the required reductions by its visit \(k=K\).\par
For any scheduling prefix, each pair’s charged work \(V_{i,j}\) is at most its largest visited \(k\), which is its number of visits. Sum over the pairs to bound total charged work by the number of visits. The all-start cap bounds actual work by charged work. The prefix through the target triple is contained in its \(h\)-sublevel, so [eq:sublevel] (\#eq:sublevel) follows. This is an inequality, not the source’s unnecessary assertion that the counts are precisely equal.\par
The non-strict biconditional in Definition 3.1, p. 14, is also too strong in the presence of ties: equal priorities in both directions force equal ranks. In the full grid, \((i,j,k)\) and \((-i,j,k)\) already tie. The paper itself describes arbitrary tie-breaking through equivalence classes on pp. 24–25. Replacing the biconditional by nondecreasing priority order makes that intended construction valid. Prefix inclusion in a sublevel is sufficient above. This transparent notation correction is not counted as another substantive finding.\par
For the polynomial full-grid criterion with exponents \(c_1,c_2>1\), put \(\tau=(|I|+1)^{c_1}(J+1)^{c_2}K\). Its sublevel count is bounded by \[2\tau\sum_{n\ge1}n^{-c_1}\sum_{n\ge1}n^{-c_2}
 \le \frac{2\tau c_1c_2}{(c_1-1)(c_2-1)}.\] Fixing a parameter removes the corresponding sum, as in Corollaries 3.5–3.6. The finite check at \(c_1=c_2=2,\tau=100\) counted 289 triples against the upper bound 800; its varying-cost trace had charged work 258 and respected the bound at every prefix. The analytic proof, not that finite enumeration, supplies the general justification.\par
\subsection*{A conservative corrected QCBP penalty}
Equation reference: eq:cone2. The following derivation supports a precise repair of F4. Define \[A_\rho=\frac{(1+\rho)^2}{1-\rho},\qquad
 B_{\rho,\gamma}=\frac{(3+\rho)\gamma}{1-\rho},\qquad
 \lambda=\frac{B_{\rho,\gamma}}{A_\rho}.\] The usual cone decomposition for \(v=z-x\), using the largest \(s\) coordinates of \(x\), gives \[\left\lVert v\right\rVert_1\le
 \frac{1+\rho}{1-\rho}
 \bigl(\left\lVert z\right\rVert_1-\left\lVert x\right\rVert_1+2\sigma_s(x)_1\bigr)
 +\frac{2\gamma\sqrt{s}}{1-\rho}\left\lVert Av\right\rVert_2.\] For completeness, the cone inequality is \(\left\lVert v_{S^c}\right\rVert_1\le\left\lVert v_S\right\rVert_1+\left\lVert z\right\rVert_1-\left\lVert x\right\rVert_1+2\sigma_s(x)_1\); combine it with the rNSP estimate \(\left\lVert v_S\right\rVert_1\le\rho\left\lVert v_{S^c}\right\rVert_1+\gamma\sqrt{s}\left\lVert Av\right\rVert_2\). Next use the largest \(s\) coordinates of \(v\), the rNSP, and \(\left\lVert v_{S^c}\right\rVert_2\le\left\lVert v\right\rVert_1/\sqrt{s}\), to obtain \[\left\lVert v\right\rVert_2\le\frac{1+\rho}{\sqrt{s}}\left\lVert v\right\rVert_1+\gamma\left\lVert Av\right\rVert_2.\] Substitution yields \[
 \left\lVert z-x\right\rVert_2\le
 \frac{A_\rho}{\sqrt{s}}
 \bigl(\left\lVert z\right\rVert_1-\left\lVert x\right\rVert_1+2\sigma_s(x)_1\bigr)
 +B_{\rho,\gamma}\left\lVert A(z-x)\right\rVert_2.\]\par
Let \(x\) be a feasible ground truth and \(\widehat x\) a QCBP minimizer, and write \(q(z)=\max\{\left\lVert Az-y\right\rVert_2-\zeta,0\}\). Then \(\left\lVert A(z-x)\right\rVert_2\le q(z)+2\zeta\) and \(\left\lVert A(\widehat x-x)\right\rVert_2\le2\zeta\). Apply [eq:cone2] (\#eq:cone2) twice, use the triangle inequality and \(\widehat f=\left\lVert \widehat x\right\rVert_1\le\left\lVert x\right\rVert_1\), and obtain \[d(z,\widehat X)\le
 \frac{A_\rho}{\sqrt{s}}
 \bigl(\left\lVert z\right\rVert_1-\widehat f+4\sigma_s(x)_1\bigr)
 +B_{\rho,\gamma}q(z)+4B_{\rho,\gamma}\zeta .\] It follows that approximate sharpness holds with \[g_Q(z)=\lambda\sqrt{s}\,q(z),\qquad
 \alpha=\frac{\sqrt{s}}{A_\rho},\quad \beta=1,\qquad
 \eta=4\sigma_s(x)_1+4\lambda\zeta\sqrt{s}.\] A larger residual coefficient also works. These constants are conservative; no claim of sharp constants is needed. The essential point is their dependence on \(\gamma\), including the coefficient in \(g_Q\). Choosing that coefficient is an application-level requirement that the original restarter’s search over \(\alpha,\beta\) alone does not automatically resolve.\par
\subsection*{Known-parameter induction, grid boundaries and rates}
For Algorithm 1, if \(\mathcal E(x_k)\le\varepsilon_k\) and \(\varepsilon_k\ge\eta\), sharpness gives \[d(x_k,\widehat X)\le
 \left(\frac{\mathcal E(x_k)+\eta}{\alpha}\right)^{1/\beta}
 \le(2\varepsilon_k/\alpha)^{1/\beta}.\] The next call is valid. Its output may fall below the noise floor; the validity premise is needed at its input, not again after this final successful call. Greedy selection prevents later interleaving from increasing \(\mathcal E\). If the target is below \(\eta\), attaining \(\eta\) suffices; if \(\eta\ge\varepsilon_0\), the initial point already meets the target. These checks support the combined-error conclusion rather than an off-by-one objection.\par
The cost sum is the restart count plus a constant times \(\sum_{k=0}^{s-1}(\varepsilon_0r^k)^{d_1/\beta-d_2}\). The finite geometric identities in all three exponent regimes were checked algebraically and corroborated by exact fractional calculations. At \(\varepsilon_0=3/2,r=1/2,s=4\), the sums for exponents \(-1,0,1\) are \(10,4,45/16\).\par
For the correct unknown-parameter pair, the source’s standing condition \(\beta\ge\beta_0\ge1\), stated on p. 14, is essential. The piecewise radius on p. 16 dominates the sharpness radius on both sides of \(2\varepsilon/\alpha_*=1\); both branches agree at that boundary. The introductory wording should be read with this standing setup. No separate defect is inferred by discarding it.\par
The paper acknowledges the loss when \(\beta_*>\beta\) on a fixed grid. For \(b=1+1/\log(1/\varepsilon)\), the critical exponent tends to zero, so the denominator of a geometric sum cannot be treated as uniformly bounded away from zero. Bounding the finite sum by its number of uniformly bounded terms gives \(K=O(\log(1/\varepsilon))\) at \(d_2=d_1/\beta\), and the stated logarithmic grid overhead follows. The \(\beta=\beta_0\) case has \(J=0\), which is compatible with the upper bound. These arguments support the qualified off-grid rates; they do not turn an iteration count into a wall-clock complexity claim.\par
\section*{Concrete methods, computations and limits}
The native method analysis checks the smoothing budget and the two primal–dual step-size/budget substitutions, in addition to F3 and F4. Its exact positive examples give a smoothing error bound \(677/4774\le1/7\), and equality at the chosen \(1/2\) error budget in each of the two primal–dual substitutions. These are bounded arithmetic checks of the stated formulas, not independent proofs of all cited convergence theorems.\par
The constrained primal–dual specialization controls the combined quantity. Its arbitrary residual weight does not independently establish \(g_Q\le\varepsilon\), consistent with F5. A tentative objection based on an extended-valued \(h\) in the unconstrained case was not promoted: the global \(Q\subseteq D\) setup restricts the domain, so the proposed indicator example did not meet the premises. Missing code, ordinary notation slips and unverified conjectures are not treated as mathematical flaws.\par
The frozen repository snapshot has commit\par
`2232e4e744efb916f7c3b3dda1faa1c5a0f38acb`\par
dated 17 April 2025. The materialized tree contains 46 files totaling 10,188,235 bytes; its retained archive has 8,246,896 bytes and SHA256\par
`bdbc67ecbc6534da2a799fa92b2a94e79215548c65d20c67982968d4bcf9750e`.\par
Inspection of the selected MATLAB scripts was static only. No author program, solver dependency or large experiment was executed. The inspected implementation material concerns NESTA and primal–dual methods; it does not supply an experimental run of the printed universal fast-gradient sequence. Numerical plots are therefore reported evidence with bounded static checks, not reproduced results.\par
The strongest conclusions here are the explicit counterexamples to literal claims. They do not establish fraud, universal failure of the restart idea, or incorrectness of every application. The correct combined-error induction, complete-grid counting under the corrected work contract, and the qualified grid-rate algebra provide substantive counterevidence to such broader conclusions.\par
\subsection*{Implementation-to-proof boundaries in the acquired snapshot}
Several static observations deserve separate qualification. The Gaussian QCBP script uses sparsity 15, whereas p. 39 specifies 10. The snapshot is later than online publication and is not established as the exact producer of the figures. Its normalized Gaussian noise lies on the sphere of the chosen radius; it is not uniform in the solid ball.\par
For the two inspected TV scripts, the analysis vector has \(2R^2\) components, so the smoothing error constant in Proposition 4.7 is \(v=R^2\). The proposition’s choice is \(\mu=\varepsilon/(2R^2)\). The scripts use \(\mu=\varepsilon/R^2\) with the same iteration budget \(\lceil8R\delta/\varepsilon\rceil\). Substitution into the stated generic estimate gives \(\varepsilon/4+\varepsilon=5\varepsilon/4\), so that proof does not directly certify the requested \(\varepsilon\) accuracy for these parameters. No iterate violating the target was constructed; this is a proof-to-implementation gap, with a source-version limitation. In the sparse \(W=I\) case the corresponding smoothing parameters do agree with the proposition.\par
The SR-LASSO scripts retain dual warm starts while using the budget from Proposition 4.12, whose stated initialization is zero dual. An additional dual-distance argument or adjusted budget is needed for that transfer. In contrast, the inspected QCBP primal–dual routine includes the term \(\kappa+\left\lVert y_0\right\rVert\) in its warm-dual budget, which accommodates the relevant triangle estimate. Neither observation constitutes an executed solver failure.\par
The schedule generator selects tuples using the powered criterion \((|i|+1)^{c_1}(j+1)^{c_2}k\), but its sorting key omits the exponents. With the documented default \(c_1=c_2=2\), the tuples \((0,0,5)\) and \((1,0,2)\) have intended priorities 5 and 8, but sorting keys 5 and 4. Thus the inspected implementation can order those tuples differently from the specified criterion. This is a concrete static discrepancy, not a failure of the abstract theorem’s corrected proof or evidence that a particular published figure used those bytes.\par
Finally, some comparison plots subtract an empirical endpoint or a CVX value and floor nonpositive differences; those baselines are not certified exact optima. The iterations-to-tolerance script uses the minimum index of a Boolean comparison and has no distinct unreached-tolerance sentinel when every comparison is true. It is unestablished whether a published run encountered that boundary. The inspected Fourier implementation uses the correct squared scaling \(AA^*=(n/m)I\), despite the local normalization transcription on p. 46. These limitations are retained in the evidence records and are not counted as additional demonstrated mathematical failures.\par
\section*{Assessment}
ID | Severity | Demonstrated consequence and limit\par
F1 | Major | The arbitrary-\(S\) work bound is false; complete vertical grids already meet the proposed structural repair.\par
F2 | Major | The conditional solver contract does not bound actual work on guessed branches; an all-start cap repairs the accounting.\par
F3 | Major | The displayed query output fails the concrete accuracy promise; returning the proximal sequence is necessary, with full theorem restoration still to check.\par
F4 | Major | General rNSP does not yield the stated fixed-penalty sharpness; an rNSP-dependent penalty is proved sufficient.\par
F5 | Minor | Combined error does not imply the stated individual feasibility rate; the formal combined bound and qualitative zero-floor feasibility survive.\par
Major classifications reflect invalid central general or application guarantees and their explicit dependencies, even where a local repair is available. The notation-only tie correction is documented separately and does not increase the substantive finding count. All findings remain subject to human adjudication. There are no human review records in this initial continuation and no independent acceptance decision.\par
\section*{Supporting checks}
The frozen source and manifest are `papers/0439/0439.pdf` and `papers/0439/source\_manifest.json`; the report source is `reports/0439.tex`. All original execution records, exact calculation outputs and source-view bindings are retained with these records.\par
Current root image observations cover physical pages 1–3, 5–8, 12–13, 16, 18–19, 21, 31–32 and 38–39; native observations retain their own attribution. Historical rendering, full extraction and programmatic input/index verification are not scientific reading or actual page viewing.\par
Actual compilation, parsing and report-page inspection are recorded in `recovery\_builds/`, `recovery\_report\_pages/`, and `report\_page\_inspection.json`. Failed and successful predecessors remain distinguished. Byte identity covers the explicit retained payloads; a generated dependency list is not a complete trace of the external environment. Publication order is frozen payloads, nonrecursive support inventory, provenance, then the root delivery entry.\par
No account UI observation, backend equivalence or independent acceptance is certified.\par
1 B. Adcock, M. J. Colbrook and M. Neyra-Nesterenko, Restarts Subject to Approximate Sharpness: A Parameter-Free and Optimal Scheme For First-Order Methods, Found. Comput. Math. 26 (2026), 1137–1192. Published online 18 February 2025. doi:10.1007/s10208-024-09673-8 (https://doi.org/10.1007/s10208-024-09673-8).\par
\end{document}
